\subsection{审计偶数窗的 Fibonacci 半单理想与幂散度谱刚性}\label{subsec:conclusion-foldbin-minideal-powerdivergence-fibonacci-rigidity}
在定理 \ref{thm:fold-bin-min-degeneracy-fib-audited} 已把审计偶数窗上的最小退化度写成 Fibonacci 律、推论 \ref{cor:fold-bin-minsector-complement-maxfiber-audited} 已把其补层规模锁定到 \(F_{m+1}=D_{2m-2}\)、而命题 \ref{prop:fold-bin-phi-identifiable-power-divergence} 已给出相对均匀基线的幂散度极限闭式之后，还可把这两条链进一步压缩为下列若干条结论：对象层的最小退化层在群胚代数中切出一个纯矩阵理想，并在阿贝尔化层形成一个规范的 Fibonacci anomaly block；子覆盖规模若沿无穷偶窗持续则最小退化层的相对秩会锁定到黄金极限；与此同时，函数层的整个 \(f_\alpha\)-散度谱则退化为一条二指数刚性曲线，并由单个 \(\chi^2\)-型极限常数完全坐标化。

\begin{theorem}[审计偶数窗中最小退化层切出的 Fibonacci 半单理想]\label{thm:conclusion-foldbin-minideal-fibonacci-primitive-splitting}
\leanverified{paper\_conclusion\_foldbin\_minideal\_fibonacci\_primitive\_splitting}
对审计偶数窗
\[
m\in\{6,8,10,12\},
\]
记
\[
A_m:=\CC[\mathcal R_m^{\mathrm{bin}}],
\qquad
\mathcal S_{m,\min}:=\mathcal S_{m,d_{\min}(m)}.
\]
在命题 \ref{prop:fold-bin-groupoid-algebra-wedderburn} 的 Wedderburn 分解
\[
A_m\cong \bigoplus_{w\in X_m}M_{d_m^{\mathrm{bin}}(w)}(\CC)
\]
中，对每个 \(w\in X_m\) 记 \(z_w\) 为第 \(w\)-块的单位元，亦即对应的中心最小幂等元，并定义
\[
e_m^{\min}:=\sum_{w\in\mathcal S_{m,\min}} z_w\in Z(A_m).
\]
则有自然同构
\[
\boxed{\;
e_m^{\min}A_m
\cong
M_{F_{m/2}}(\CC)^{\oplus F_m}.\;}
\]
并且其补理想满足
\[
\boxed{\;
\dim_{\CC}Z\bigl((1-e_m^{\min})A_m\bigr)
=
F_{m+1}
=
D_{2m-2}.\;}
\]
若再记 \(\pi_w:A_m\to M_{d_m^{\mathrm{bin}}(w)}(\CC)\) 为第 \(w\)-块坐标投影，并定义
\[
\operatorname{Prim}_{\min}(A_m):=\{\ker\pi_w:\ w\in\mathcal S_{m,\min}\},
\qquad
\operatorname{Prim}_{\mathrm{rest}}(A_m):=\{\ker\pi_w:\ w\in X_m\setminus\mathcal S_{m,\min}\},
\]
则 primitive spectrum 分裂为
\[
\boxed{\;
\operatorname{Prim}(A_m)
=
\operatorname{Prim}_{\min}(A_m)\sqcup \operatorname{Prim}_{\mathrm{rest}}(A_m),\qquad
|\operatorname{Prim}_{\min}(A_m)|=F_m,\quad
|\operatorname{Prim}_{\mathrm{rest}}(A_m)|=F_{m+1}.\;}
\]
\end{theorem}
\begin{proof}
由命题 \ref{prop:fold-bin-groupoid-algebra-wedderburn}，
\[
A_m\cong \bigoplus_{w\in X_m}M_{d_m^{\mathrm{bin}}(w)}(\CC),
\]
故
\[
e_m^{\min}A_m
\cong
\bigoplus_{w\in\mathcal S_{m,\min}}M_{d_m^{\mathrm{bin}}(w)}(\CC).
\]
再由定理 \ref{thm:fold-bin-min-degeneracy-fib-audited}，
\[
d_m^{\mathrm{bin}}(w)=d_{\min}(m)=F_{m/2}
\qquad (w\in\mathcal S_{m,\min}),
\]
并且
\[
|\mathcal S_{m,\min}|=F_m.
\]
于是得到
\[
e_m^{\min}A_m\cong M_{F_{m/2}}(\CC)^{\oplus F_m}.
\]

另一方面，补理想
\[
(1-e_m^{\min})A_m
\cong
\bigoplus_{w\in X_m\setminus\mathcal S_{m,\min}}M_{d_m^{\mathrm{bin}}(w)}(\CC)
\]
的中心维数等于其 Wedderburn 块数，即
\[
|X_m\setminus\mathcal S_{m,\min}|.
\]
由推论 \ref{cor:fold-bin-minsector-complement-maxfiber-audited}，
\[
|X_m\setminus\mathcal S_{m,\min}|=F_{m+1}=D_{2m-2},
\]
故得到第二条盒装公式。

最后，有限维半单复代数的 primitive ideals 恰由各矩阵块的坐标核给出，因此
\[
\operatorname{Prim}(A_m)=\{\ker\pi_w:\ w\in X_m\}.
\]
把这一集合按 \(w\in\mathcal S_{m,\min}\) 与 \(w\notin\mathcal S_{m,\min}\) 分拆，即得
\[
\operatorname{Prim}(A_m)=\operatorname{Prim}_{\min}(A_m)\sqcup \operatorname{Prim}_{\mathrm{rest}}(A_m),
\]
并且两部分基数分别等于对应块数 \(F_m\) 与 \(F_{m+1}\)。证毕。
\end{proof}

\begin{proposition}[已审计偶窗中的最小退化扇区给出规范 Fibonacci anomaly block]\label{prop:conclusion-audited-even-window-fibonacci-anomaly-block}
\leanverified{paper\_conclusion\_audited\_even\_window\_fibonacci\_anomaly\_block}
对已审计偶数窗
\[
m\in\{6,8,10,12\},
\]
定义
\[
\mathfrak A_m^{\min}
:=
\bigoplus_{w\in \mathcal S_{m,\min}}\FF_2 e_w
\subset
\bigl(G_m^{\mathrm{bin}}\bigr)^{\mathrm{ab}},
\]
其中
\[
\mathcal S_{m,\min}:=\mathcal S_{m,d_{\min}(m)}
\]
为最小退化扇区，而 \(e_w\) 表示阿贝尔化中的第 \(w\)-坐标基向量。则 \(\mathfrak A_m^{\min}\) 是一个规范坐标子晶格，并满足
\[
\dim_{\FF_2}\mathfrak A_m^{\min}=F_m.
\]
更具体地，
\[
\dim_{\FF_2}\mathfrak A_6^{\min}=8,\qquad
\dim_{\FF_2}\mathfrak A_8^{\min}=21,\qquad
\dim_{\FF_2}\mathfrak A_{10}^{\min}=55,\qquad
\dim_{\FF_2}\mathfrak A_{12}^{\min}=144.
\]
该子晶格完全由最小退化纤维决定，而与较高退化扇区无关。
\end{proposition}
\begin{proof}
由定理 \ref{thm:fold-bin-min-degeneracy-fib-audited}，
\[
d_{\min}(m)=F_{m/2},
\qquad
|\mathcal S_{m,\min}|=F_m
\]
对 \(m\in\{6,8,10,12\}\) 成立。特别地，
\[
d_{\min}(m)\ge 2.
\]
因此在这些审计偶窗内，一切纤维都满足 \(d_m^{\mathrm{bin}}(w)\ge 2\)。由 bin-fold gauge 群阿贝尔化的坐标描述，
\[
\bigl(G_m^{\mathrm{bin}}\bigr)^{\mathrm{ab}}
\cong
(\ZZ/2\ZZ)^{|X_m|}
\cong
\bigoplus_{w\in X_m}\FF_2 e_w.
\]
于是取最小退化扇区对应的坐标子空间
\[
\mathfrak A_m^{\min}
=
\bigoplus_{w\in \mathcal S_{m,\min}}\FF_2 e_w
\]
即可得到一个规范定义的 anomaly block，其维数恰为
\[
\dim_{\FF_2}\mathfrak A_m^{\min}=|\mathcal S_{m,\min}|=F_m.
\]
将 \(m=6,8,10,12\) 代入即得四个显式数值。证毕。
\end{proof}

\begin{corollary}[正则子覆盖外推下最小 anomaly block 的相对秩极限]\label{cor:conclusion-minimal-anomaly-block-relative-rank-phi-square-inverse}
\leanverified{paper\_conclusion\_minimal\_anomaly\_block\_relative\_rank\_phi\_square\_inverse}
若最小退化扇区的正则子覆盖外推在一列趋于无穷的偶整数 \(m\) 上成立，从而
\[
\frac{|\mathcal S_{m,\min}|}{|X_m|}\longrightarrow \varphi^{-2},
\]
则沿该子列有
\[
\frac{\dim_{\FF_2}\mathfrak A_m^{\min}}
{\dim_{\FF_2}\bigl(G_m^{\mathrm{bin}}\bigr)^{\mathrm{ab}}}
\longrightarrow
\varphi^{-2}.
\]
\end{corollary}
\begin{proof}
由命题 \ref{prop:conclusion-audited-even-window-fibonacci-anomaly-block}，
\[
\dim_{\FF_2}\mathfrak A_m^{\min}=|\mathcal S_{m,\min}|.
\]
另一方面，在所考虑的偶窗上，由 \(d_{\min}(m)\ge 2\) 可知所有纤维都为阿贝尔化贡献一个 \(\FF_2\)-坐标，因此
\[
\dim_{\FF_2}\bigl(G_m^{\mathrm{bin}}\bigr)^{\mathrm{ab}}=|X_m|.
\]
两式相除，再用假设中的极限
\[
\frac{|\mathcal S_{m,\min}|}{|X_m|}\to \varphi^{-2},
\]
便得到结论。证毕。
\end{proof}

\begin{theorem}[子覆盖规模持续时最小退化度相对于 Fibonacci 候选律的指数塌缩]\label{thm:conclusion-foldbin-subcover-exponential-mincollapse}
\leanverified{paper\_conclusion\_foldbin\_subcover\_exponential\_mincollapse}
设沿某个趋于无穷的偶数子序列 \(m\to\infty\) 仍有
\[
\bigl|\mathcal S_{m,d_{\min}(m)}\bigr|
=
|X_{m-2}|
=
F_m.
\]
则必有
\[
\boxed{\;
\frac{d_{\min}(m)}{F_{m/2}}
\le
\frac{2^m}{F_mF_{m/2}}
=
O\!\left(\left(\frac{2}{\varphi^{3/2}}\right)^m\right)
\longrightarrow 0.\;}
\]
特别地，只要最小退化扇区的块数沿无穷偶窗保持子覆盖规模 \(F_m\)，最小退化度就不可能与候选 Fibonacci 律 \(F_{m/2}\) 同阶，更不可能持续取等。
\end{theorem}
\begin{proof}
由总预像计数恒等式
\[
2^m=\sum_{w\in X_m}d_m^{\mathrm{bin}}(w)
\]
与最小退化度定义，得到
\[
2^m\ge \bigl|\mathcal S_{m,d_{\min}(m)}\bigr|\,d_{\min}(m).
\]
在假设下，
\[
\bigl|\mathcal S_{m,d_{\min}(m)}\bigr|=F_m,
\]
故
\[
d_{\min}(m)\le \frac{2^m}{F_m}.
\]
再除以 \(F_{m/2}\)，便有
\[
\frac{d_{\min}(m)}{F_{m/2}}
\le
\frac{2^m}{F_mF_{m/2}}.
\]
由 Binet 渐近
\[
F_n\asymp \frac{\varphi^n}{\sqrt5},
\]
可得
\[
F_mF_{m/2}\asymp \varphi^{3m/2},
\]
从而
\[
\frac{2^m}{F_mF_{m/2}}
=
O\!\left(\left(\frac{2}{\varphi^{3/2}}\right)^m\right).
\]
又因
\[
\varphi^3=2+\sqrt5>4,
\]
故 \(\varphi^{3/2}>2\)，于是右端指数趋于 \(0\)。证毕。
\end{proof}

\begin{theorem}[\texorpdfstring{$f_\alpha$}{f_alpha} 散度极限常数的二阶线性递推与严格对数凸性]\label{thm:conclusion-power-divergence-secondorder-recurrence}
对 \(\alpha>1\) 定义
\[
f_\alpha(t):=t^\alpha-1-\alpha(t-1),
\qquad
E_\alpha:=D_{f_\alpha}^\infty+1.
\]
\leanverified{paper\_conclusion\_power\_divergence\_secondorder\_recurrence}
则有显式表示
\[
E_\alpha
=
\varphi^{-1}\left(\frac{\varphi^2}{\sqrt5}\right)^\alpha
+
\varphi^{-2}\left(\frac{\varphi}{\sqrt5}\right)^\alpha.
\]
从而对一切 \(\alpha>1\) 都成立精确递推
\[
\boxed{\;
E_{\alpha+2}
=
\left(1+\frac{2}{\sqrt5}\right)E_{\alpha+1}
-\frac{2+\sqrt5}{5}\,E_\alpha.\;}
\]
并且对任意 \(\alpha\neq \beta\) 与 \(\theta\in(0,1)\)，有严格对数凸不等式
\[
\boxed{\;
E_{\theta\alpha+(1-\theta)\beta}
<
E_\alpha^\theta E_\beta^{1-\theta}.\;}
\]
特别地，\(\alpha\mapsto E_\alpha\) 落在二维线性张成
\[
\operatorname{span}\left\{
\left(\frac{\varphi^2}{\sqrt5}\right)^\alpha,\,
\left(\frac{\varphi}{\sqrt5}\right)^\alpha
\right\}
\]
之中；因此任意两点
\[
(\alpha_0,E_{\alpha_0}),\qquad (\alpha_1,E_{\alpha_1})
\qquad (\alpha_0\neq \alpha_1)
\]
唯一决定整条 \(E\)-谱，等价地唯一决定整条幂散度极限曲线。
\end{theorem}
\begin{proof}
由命题 \ref{prop:fold-bin-phi-identifiable-power-divergence}，
\[
D_{f_\alpha}^\infty
=
\frac{1}{5^{\alpha/2}}\bigl(\varphi^{2\alpha-1}+\varphi^{\alpha-2}\bigr)-1,
\]
故
\[
E_\alpha
=
\varphi^{-1}\left(\frac{\varphi^2}{\sqrt5}\right)^\alpha
+
\varphi^{-2}\left(\frac{\varphi}{\sqrt5}\right)^\alpha.
\]
记
\[
r:=\frac{\varphi^2}{\sqrt5},
\qquad
s:=\frac{\varphi}{\sqrt5},
\qquad
A:=\varphi^{-1},
\qquad
B:=\varphi^{-2},
\]
则
\[
E_\alpha=Ar^\alpha+Bs^\alpha.
\]
因此
\[
E_{\alpha+2}=(r+s)E_{\alpha+1}-rs\,E_\alpha.
\]
直接计算
\[
r+s=\frac{\varphi^2+\varphi}{\sqrt5}
=
\frac{2+\sqrt5}{\sqrt5}
=
1+\frac{2}{\sqrt5},
\]
\[
rs=\frac{\varphi^3}{5}
=
\frac{2+\sqrt5}{5},
\]
便得到所示二阶线性递推。

再把 \(E_\alpha\) 写成
\[
E_\alpha=Ae^{c\alpha}+Be^{d\alpha},
\qquad
c:=\log r,\quad d:=\log s.
\]
由于 \(A,B>0\) 且 \(c\neq d\)，有
\[
\frac{d^2}{d\alpha^2}\log E_\alpha
=
\frac{AB(c-d)^2e^{(c+d)\alpha}}{(Ae^{c\alpha}+Be^{d\alpha})^2}
>0.
\]
故 \(\alpha\mapsto \log E_\alpha\) 严格凸，从而
\[
E_{\theta\alpha+(1-\theta)\beta}<E_\alpha^\theta E_\beta^{1-\theta}
\qquad
(\alpha\neq\beta,\ \theta\in(0,1)).
\]

最后，对任意 \(\alpha_0\neq \alpha_1\)，由
\[
\begin{pmatrix}
E_{\alpha_0}\\[2pt]
E_{\alpha_1}
\end{pmatrix}
=
\begin{pmatrix}
r^{\alpha_0} & s^{\alpha_0}\\
r^{\alpha_1} & s^{\alpha_1}
\end{pmatrix}
\begin{pmatrix}
A\\
B
\end{pmatrix}
\]
可知只要上式中的 \(2\times 2\) 矩阵可逆，\((A,B)\) 就由两点值唯一决定。而其行列式为
\[
r^{\alpha_0}s^{\alpha_1}-r^{\alpha_1}s^{\alpha_0}
=(rs)^{\alpha_0}\bigl(s^{\alpha_1-\alpha_0}-r^{\alpha_1-\alpha_0}\bigr)\neq 0,
\]
因为 \(r\neq s\) 且 \(\alpha_0\neq \alpha_1\)。于是整条 \(E\)-谱被唯一固定，而 \(D_{f_\alpha}^\infty=E_\alpha-1\) 也随之唯一固定。证毕。
\end{proof}

\begin{corollary}[单个 \texorpdfstring{$\chi^2$}{chi2} 极限常数对黄金参数与整条幂散度谱的恢复]\label{cor:conclusion-chi2-recovers-full-power-divergence-family}
取 \(\alpha=2\)。则
\[
\boxed{\;
D_{f_2}^\infty
=
\frac{\sqrt5-2}{5},
\qquad
\sqrt5=5D_{f_2}^\infty+2,
\qquad
\varphi=\frac{3+5D_{f_2}^\infty}{2}.\;}
\]
因此对任意 \(\alpha>1\)，都有
\[
\boxed{\;
D_{f_\alpha}^\infty
=
\frac{
\left(\frac{3+5D_{f_2}^\infty}{2}\right)^{2\alpha-1}
+
\left(\frac{3+5D_{f_2}^\infty}{2}\right)^{\alpha-2}
}{
\left(5D_{f_2}^\infty+2\right)^\alpha
}
-1.\;}
\]
换言之，相对均匀基线的 \(\chi^2\)-极限常数单独就足以代数地恢复黄金参数 \(\varphi\)，并进而恢复全部幂散度极限。
\end{corollary}
\begin{proof}
由定理 \ref{thm:conclusion-power-divergence-secondorder-recurrence} 的显式公式，
\[
D_{f_2}^\infty
=
\frac{1}{5}\bigl(\varphi^3+1\bigr)-1.
\]
再用
\[
\varphi^2=\varphi+1,
\qquad
\varphi^3=\varphi(\varphi+1)=2\varphi+1=2+\sqrt5,
\]
便得到
\[
D_{f_2}^\infty=\frac{\sqrt5-2}{5}.
\]
移项即得
\[
\sqrt5=5D_{f_2}^\infty+2,
\qquad
\varphi=\frac{1+\sqrt5}{2}=\frac{3+5D_{f_2}^\infty}{2}.
\]
最后把这两个表达式代回
\[
D_{f_\alpha}^\infty
=
\frac{1}{5^{\alpha/2}}\bigl(\varphi^{2\alpha-1}+\varphi^{\alpha-2}\bigr)-1
\]
即可得到所示闭式。证毕。
\end{proof}

\begin{theorem}[审计偶窗首折点的同伦纯化]\label{thm:conclusion-audited-even-firstkink-homotopy-purification}
\leanverified{paper\_conclusion\_audited\_even\_firstkink\_homotopy\_purification}
对审计偶窗
\[
m\in\{6,8,10,12\},
\]
记
\[
d_{\min}(m)=F_{m/2},
\qquad
\mathcal S_{m,\min}:=\mathcal S_{m,d_{\min}(m)},
\]
并定义
\[
C_m(B):=C_m^{\mathrm{cont}}(2^B),
\qquad
\mathrm{Succ}^{\mathrm{bin}}_m(B):=2^{-m}C_m(B).
\]
则当
\[
0\le B\le \log_2 d_{\min}(m)=\log_2 F_{m/2}
\]
时，有精确公式
\[
C_m(B)=2^B F_{m+2},
\qquad
\mathrm{Succ}^{\mathrm{bin}}_m(B)=2^{B-m}F_{m+2}.
\]
同时，
\[
|\mathcal S_{m,\min}|=F_m=|X_{m-2}|,
\]
并且对任意 \(x\in \mathcal S_{m,\min}\)，纤维独立集复形满足
\[
|K_x|\simeq *
\qquad (m\equiv 0\!\!\!\pmod 6),
\]
\[
|K_x|\simeq S^{\tau(x)-1}
\qquad (m\equiv 2,4\!\!\!\pmod 6).
\]
因而在这些审计偶窗内，首个容量折点之前的线性容量段与最小退化扇区的纯同伦相位在同一 Fibonacci 界面上同步出现。
\end{theorem}
\begin{proof}
由容量公式
\[
C_m(B)=\sum_{w\in X_m}\min\{d_m^{\mathrm{bin}}(w),2^B\}
\]
与定理 \ref{thm:fold-bin-min-degeneracy-fib-audited}，
\[
d_m^{\mathrm{bin}}(w)\ge d_{\min}(m)=F_{m/2}
\qquad (w\in X_m).
\]
若 \(0\le B\le \log_2 d_{\min}(m)\)，则 \(2^B\le d_{\min}(m)\le d_m^{\mathrm{bin}}(w)\) 对全部 \(w\in X_m\) 成立，从而每一项都等于 \(2^B\)。因此
\[
C_m(B)=2^B|X_m|=2^BF_{m+2}.
\]
再除以 \(2^m\) 即得
\[
\mathrm{Succ}^{\mathrm{bin}}_m(B)=2^{-m}C_m(B)=2^{B-m}F_{m+2}.
\]

另一方面，定理 \ref{thm:fold-bin-min-degeneracy-fib-audited} 已给出
\[
|\mathcal S_{m,\min}|=F_m=|X_{m-2}|,
\]
而推论 \ref{cor:xi-time-part55ab-audited-even-minsector-strict-homotopy-transition} 给出
\[
|K_x|\simeq *
\quad (m\equiv0\!\!\!\pmod6),
\qquad
|K_x|\simeq S^{\tau(x)-1}
\quad (m\equiv2,4\!\!\!\pmod6)
\]
对所有 \(x\in\mathcal S_{m,\min}\) 同时成立。合并上述两部分即得结论。证毕。
\end{proof}


\begin{theorem}[完整异常账本与近无损逆码预算的指数分离]\label{thm:conclusion-binfold-anomaly-ledger-vs-nearlossless-budget-exponential-separation}
对每个分辨率 \(m\)，定义
\[
N_m:=\#\{w\in X_m:\ d_m(w)\ge 2\}.
\]
若满足
\[
d_{\min}(m):=\min_{w\in X_m}d_m(w)\ge 2,
\]
则
\[
N_m=|X_m|,
\qquad
\bigl(G_m^{\mathrm{bin}}\bigr)^{\mathrm{ab}}
\cong
(\ZZ/2\ZZ)^{N_m}.
\]
另一方面，若 \(B_m(\varepsilon)\) 是任意满足
\[
\mathrm{Succ}_m\!\bigl(B_m(\varepsilon)\bigr)\ge 1-\varepsilon
\qquad (0<\varepsilon<1)
\]
的可逆侧信息预算，则必有
\[
B_m(\varepsilon)\ge m-\log_2|X_m|+\log_2(1-\varepsilon).
\]
因此
\[
\frac{\dim_{\FF_2}\bigl(G_m^{\mathrm{bin}}\bigr)^{\mathrm{ab}}}{B_m(\varepsilon)}
\ge
\frac{|X_m|}{m-\log_2|X_m|+O(1)}.
\]
特别地，沿任意满足 \(d_{\min}(m)\ge 2\) 的无穷子序列，若
\[
|X_m|=F_{m+2}\sim \frac{\varphi^m}{\sqrt5},
\]
则
\[
\frac{\dim_{\FF_2}\bigl(G_m^{\mathrm{bin}}\bigr)^{\mathrm{ab}}}{B_m(\varepsilon)}
\]
按指数速度发散。换言之，完整异常账本的最小完备维数与近无损微态逆码的必要预算之间存在严格的指数级资源分离。
\end{theorem}
\begin{proof}
若 \(d_{\min}(m)\ge 2\)，则每个稳定类型 \(w\in X_m\) 都满足 \(d_m(w)\ge 2\)，从而
\[
N_m=\#\{w\in X_m:\ d_m(w)\ge 2\}=|X_m|.
\]
与命题 \ref{prop:conclusion-audited-even-window-fibonacci-anomaly-block} 的证明完全相同，bin-fold gauge 群阿贝尔化的坐标描述给出
\[
\bigl(G_m^{\mathrm{bin}}\bigr)^{\mathrm{ab}}
\cong
(\ZZ/2\ZZ)^{|X_m|}
=
(\ZZ/2\ZZ)^{N_m},
\]
故
\[
\dim_{\FF_2}\bigl(G_m^{\mathrm{bin}}\bigr)^{\mathrm{ab}}=N_m=|X_m|.
\]

另一方面，定理 \ref{thm:conclusion-tail-capacity-moment-tail-law-equivalence} 已证明：一旦
\[
\mathrm{Succ}_m(B)\ge 1-\varepsilon,
\]
就必有
\[
B\ge m-\log_2|X_m|+\log_2(1-\varepsilon).
\]
将 \(B=B_m(\varepsilon)\) 代入，并利用上面的维数公式，便得到
\[
\frac{\dim_{\FF_2}\bigl(G_m^{\mathrm{bin}}\bigr)^{\mathrm{ab}}}{B_m(\varepsilon)}
\ge
\frac{|X_m|}{m-\log_2|X_m|+\log_2(1-\varepsilon)}.
\]
由于 \(\varepsilon\) 固定，分母可写成
\[
m-\log_2|X_m|+O(1),
\]
故得到所示不等式。

最后，若沿某无穷子序列有
\[
|X_m|=F_{m+2}\sim \frac{\varphi^m}{\sqrt5},
\]
则分子按 \(\varphi^m\) 指数增长，而分母至多线性于 \(m\)。因此该比值必按指数速度发散。证毕。
\end{proof}

\noindent
因此，在 bin-fold 口径下出现了四类彼此咬合的 Fibonacci 刚性。其一，审计偶数窗的最小退化层在对象层切出一个由 \(F_{m/2}\) 与 \(F_m\) 精确控制的纯矩阵理想，并在阿贝尔化层诱导出一个规范的 Fibonacci anomaly block；其相对秩在正则子覆盖外推下进一步锁定到 \(\varphi^{-2}\)。其二，首个容量折点之前的连续容量函数严格保持斜率 \(2^B\mapsto 2^BF_{m+2}\)，而最小退化扇区又在同一窗口内发生模 \(6\) 纯同伦相位切换，从而“首折点”与“纯同伦界面”被压到同一最小扇区上。其三，相对均匀基线的幂散度极限谱在函数层退化为二指数系统，既满足精确二阶递推，又能由单个 \(\chi^2\)-型常数恢复 \(\varphi\) 并回收整条谱。其四，一旦全部纤维都至少带有二重退化，则完整异常账本的最小完备维数会随 \(|X_m|\) 以 Fibonacci 指数膨胀，而近无损逆码预算仍然只受 \(m-\log_2|X_m|\) 级的必要下界控制，于是二者之间出现严格的指数级资源分离。由此，最小退化层的对象刚性、容量首折点的计数刚性、输出非均匀性的函数刚性与异常账本/逆码预算的资源刚性，便在同一 Fibonacci 机制下合并为彼此对应的四侧表述。

\endinput
